\subsection{规则的描述难度与计算难度}

\subsubsection{两个不同的困难来源}
一个任务可以容易形式化，却需要很大的计算量；也可以由人较容易完成，却难以写出完整规则。国际象棋的棋盘、棋子与合法走法都很明确，其困难主要来自搜索规模。人脸检测与语音识别则说明了另一种困难：识别结果容易理解，但把所有输入条件转化为固定规则很困难。文学分析与国际外交被用作更复杂任务的示意例子。

这种区分与并行计算有关。对于运算规则已经明确的任务，可以继续研究算法复杂度、并行性与数据带宽；对于输入输出关系本身难以手工表达的任务，还需要解决模型如何从观测中获得这种关系的问题。

\subsubsection{国际象棋：规则明确，搜索树很大}
棋盘有 64 个位置，初始有 32 个棋子。根据合法走法生成后继局面、为搜索树叶节点评分、再根据这些评分选择路径，就形成了一种可明确描述的计算过程。如果每个局面平均约有 30 种合法走法，搜索 10 个半回合的叶节点数量约为
\begin{equation}
30^{10}=5.9049\times10^{14}.
\end{equation}
即使每秒处理 $10^6$ 个局面，也需要约 $5.9\times10^8$ 秒，即约 18.7 年。因此，启发式方法、剪枝、并行搜索和更快的计算机共同影响这种方案是否实用。Deep Blue 在 1997 年击败 Garry Kasparov，是一个历史例子。

\subsubsection{规则系统：常识需要一致的语义表示}
规则系统用事实与规则驱动推理引擎（inference engine），产生可由已有知识推出的结论。Cyc 的例子把常识组织为本体和知识库，再对形式化陈述进行推理。例如，从“某人是一位美国总统”和“美国总统是美国公民”可以推出该人是美国公民。这说明规则能够显式表达一条推理链。

“Fred 正在刮胡子”的例子揭示了语义表示的难处：人没有电气部件，人刮胡子时有剃刀，剃刀有电气部件。若把“拥有某物”与“某物是身体的组成部分”混为同一种关系，规则之间就可能产生错误关联。能够处理真实常识的系统需要大量事实、规则及一致的关系定义；这里的规则系统依赖预先给定的规则与陈述。

\subsection{人工特征与学习得到的表示}

\subsubsection{经典机器学习中的特征映射}
经典机器学习通常先由人选定特征，再学习特征与输出之间的映射。例如，朴素贝叶斯（naïve Bayes）采用“给定类别后，各特征条件独立”的假设；逻辑回归（logistic regression）学习各特征对输出的权重，训练过程可以使用梯度下降。这两个例子说明，特征如何表示输入，与模型如何利用这些特征，是可以分别考虑的环节。

不同任务需要不同的特征。图像可以经过视觉特征提取再进行检测；音频可以转换成音频特征，用于识别说话人；文本特征可以用于文本分类、机器翻译和信息检索。特征必须保留与任务输出有关的信息，同时让后续映射更容易表达。

\subsubsection{坐标表示能够改变分类边界的复杂度}
图~\ref{l10:fig:representation} 中，蓝色点聚集在中心，绿色点分布在外围。在直角坐标 $(x,y)$ 下，仅用一条直线难以将两者分开；图中的多条线性不等式是在尝试用若干半平面描述中心区域。

\begin{figure}[H]
\centering
\begin{subfigure}[t]{0.48\linewidth}
\centering\includegraphics[height=2.18in,width=\linewidth,keepaspectratio]{representation_xy.png}
\caption{直角坐标表示。}
\end{subfigure}\hfill
\begin{subfigure}[t]{0.48\linewidth}
\centering\includegraphics[height=2.18in,width=\linewidth,keepaspectratio]{representation_polar.png}
\caption{半径与角度表示。}
\end{subfigure}
\caption{相同分类任务在不同特征坐标下具有不同的分离难度。}
\label{l10:fig:representation}
\end{figure}

若改用半径 $r=\sqrt{x^2+y^2}$ 和角度 $\theta$，中心点与外围点的差异就主要表现为 $r$ 的大小。于是，用一个半径阈值就能描述图示的分离关系。这里有用的特征是“到中心的距离”，它把几何结构直接交给后续分类器。

图中的 $\theta=\arctan(y/x)$ 是简写；完整角度表示还需要区分象限，并处理 $x=0$ 的情形。该例用于说明半径特征的作用，不依赖具体的角度实现。

\subsubsection{从变化因素到层次化表示}
人工指定“有车轮”作为识别汽车的特征，仍然要回答如何从像素中识别车轮。圆形轮廓、较暗的边缘等规则又会受到遮挡、观察角度、阴影和轮胎外观的影响。因此，选定一个看似合理的语义特征，并没有自动解决其底层检测问题。

变化因素（factors of variation）指影响观测数据的潜在因素。例如，年龄、性别、口音和说出的词语会影响语音；位置、颜色和太阳角度会影响汽车图像。一个观测通常同时受到多个因素影响，学习表示需要从这些混合变化中提取与任务相关的结构。

表示学习（representation learning）把特征本身也纳入学习过程。深度学习（deep learning）进一步用其他表示来表达新的表示，形成层次化特征。图像中的示意层次可以是输入像素、边缘、角点与轮廓、对象部件，最后是对象类别；较复杂概念由较简单概念组合而成。

\begin{figure}[H]
\centering
\includegraphics[width=0.52\linewidth]{ml_evolution.png}
\caption{规则系统、经典机器学习与深度学习中的输入到输出路径。深度学习在中间引入多层特征表示。}
\label{l10:fig:evolution}
\end{figure}

\begin{keybox}[title={输入输出关系与中间表示}]
规则系统直接编码输入输出关系；经典机器学习学习人工特征到输出的映射；表示学习还学习特征的构造方式。深度学习通过层次组合表达特征。多层感知机中的隐藏层，为这种“由输入构造新表示，再根据新表示做判断”的过程提供了一个具体模型。
\end{keybox}
